%%%%%%%%%%%%%%%%%%%%%%%%%%%%%%%%%%%%%%%%%%%%%%%%%%%%%%%%%%%%
\documentclass[xcolor=x11names,compress]{beamer}

\definecolor{CoolBlack}{rgb}{0.0, 0.18, 0.39}
\definecolor{byellow}{rgb}{0.55037, 0.38821, 0.06142}
%% General document %%%%%%%%%%%%%%%%%%%%%%%%%%%%%%%%%%
\usepackage{graphicx}
\usepackage{tikz}
\usepackage{Tabbing}
\usetikzlibrary{decorations.fractals}
%%%%%%%%%%%%%%%%%%%%%%%%%%%%%%%%%%%%%%%%%%%%%%%%%%%%%%

%% Beamer Layout %%%%%%%%%%%%%%%%%%%%%%%%%%%%%%%%%%
\useoutertheme[subsection=false,shadow]{miniframes}
\useinnertheme{default}
\usefonttheme{serif}
\usepackage{palatino}
\usepackage{tabu}

% addition of color
\usepackage{xcolor}
\definecolor{dgreen}{rgb}{0.,0.6,0.}
\definecolor{RawSienna}{cmyk}{0,0.72,1,0.45}

\setbeamerfont{title like}{shape=\scshape}
\setbeamerfont{frametitle}{shape=\scshape}

\setbeamercolor*{lower separation line head}{bg=CoolBlack} 
\setbeamercolor*{normal text}{fg=black,bg=white} 
\setbeamercolor*{alerted text}{fg=red} 
\setbeamercolor*{example text}{fg=black} 
\setbeamercolor*{structure}{fg=black} 
 
\setbeamercolor*{palette tertiary}{fg=black,bg=black!10} 
\setbeamercolor*{palette quaternary}{fg=black,bg=black!10} 

\renewcommand{\(}{\begin{columns}}
\renewcommand{\)}{\end{columns}}
\newcommand{\<}[1]{\begin{column}{#1}}
\renewcommand{\>}{\end{column}}

% adding slide numbers
\addtobeamertemplate{navigation symbols}{}{%
    \usebeamerfont{footline}%
    \usebeamercolor[fg]{footline}%
    \hspace{1em}%
    \insertframenumber/\inserttotalframenumber
}

% equation stuff
\newcommand{\Macro}{\ensuremath{\Sigma}}
\newcommand{\Sn}{\ensuremath{S_N} }
\newcommand{\vOmega}{\ensuremath{\hat{\Omega}}}
\usepackage{mathrsfs}
\usepackage[mathcal]{euscript}
\usepackage{amssymb}
\usepackage{amsthm}
\usepackage{epsfig}
\usepackage{amsmath}

\newcommand{\ve}[1]{\ensuremath{\mathbf{#1}}}
\newcommand{\micro}{\ensuremath{\sigma}}
\newcommand{\detR}{\ensuremath{\Sigma}}
%%%%%%%%%%%%%%%%%%%%%%%%%%%%%%%%%%%%%%%%%%%%%%%%%%

\begin{document}

%%%%%%%%%%%%%%%%%%%%%%%%%%%%%%%%%%%%%%%%%%%%%%%%%%%%%%
%%%%%%%%%%%%%%%%%%%%%%%%%%%%%%%%%%%%%%%%%%%%%%%%%%%%%%
\begin{frame}
\title{Monte Carlo Variance Reduction Intro}
\subtitle{NE 155, Spring 2014}
\author{
        \includegraphics[height=2cm]{bk}\\R.\ N.\ Slaybaugh}

\date{April 30, 2014}
\titlepage
\end{frame}


% --------------------------------------------------------------
% --------------------------------------------------------------
\section{\scshape Variance Reduction}
\begin{frame}[fragile]
  \frametitle{Acceleration}
  \begin{itemize}
  	\item To use MC in many applications, we need to \textit{accelerate} it
	\item Variance reduction is designed to improve the FOM:
  \end{itemize}
  \begin{align}
  \text{FOM} = \frac{1}{\text{R}^2\text{t}} \qquad & \text{R = relative error} \nonumber \\ 
  & \text{t = time} \nonumber 
  \end{align}
  \begin{itemize}
  	\item \underline{Goal:} improve FOM
  	\item Control how and when particles contribute to tallies in a way that does not bias physics (``fair game") 
  \end{itemize}

\end{frame}


% --------------------------------------------------------------
% --------------------------------------------------------------
\begin{frame}[fragile]
  \frametitle{Influence of Weight}
  \begin{itemize}
  	\item Weight corresponds to influence on answer
  \end{itemize}
  \begin{align}
  \phi_c V = \frac{\sum_{i=1}^{n_c} w_{i} }{N \Sigma_t} \nonumber 
  \end{align}
  \begin{itemize}
  	\item a particle with \textbf{high} weight contributes a lot
  	\\ \underline{concern?} can increase variance
  	\item a particle with \textbf{low} weight contributes very little
  	\\ \underline{concern?} can waste time  	
  	\item Variance Reduction: track particles that will contribute meaningfully; discard those that won't; maintain a fair game
  \end{itemize}

\end{frame}


% --------------------------------------------------------------
% --------------------------------------------------------------
\begin{frame}[fragile]
  \frametitle{Survival Biasing}
  \begin{itemize}
  	\item Probability of absorption at a collision?
  \end{itemize}

\begin{tabu}{|X | c |X|}
\hline
 & Analog &  Survival Biasing \\\hline
Tally at collision & $w_i$ &  \[w_i \frac{\Sigma_a}{\Sigma_t}\]  \\
$w_{i+1}$ & $0$ & \[w_i\bigl(1-\frac{\Sigma_a}{\Sigma_t} \bigr)\] \\\hline
\end{tabu}

  \begin{itemize}
  	\item maintains a fair game
  	\item represents physics
  	\item more information per history
  \end{itemize}

\end{frame}


% --------------------------------------------------------------
% --------------------------------------------------------------
\begin{frame}[fragile]
  \frametitle{Splitting}
  
  If a particle has too high of a weight, split it.
  \vspace*{1 em}
  
  \begin{itemize}
  	\item if $w_i > w_{\max}$, split:
  	\begin{itemize}
  	    	\item Split Ratio, SR, $\equiv \frac{w_i}{w_{nom}}$
  	\item get $\xi$, a random number on $0 \leq \xi \leq 1$
  	\item if $\xi \geq$ (SR - int(SR)):\\
  	      \hspace*{2 em} create int(SR) new particles 
  	\item else:\\
  	      \hspace*{2 em}create in(SR) + 1 new particles
  	\item $\forall$ new particles, $w_{i+1} = \frac{w_i}{\text{\# new particles}}$
  	\end{itemize}
  \end{itemize}

\end{frame}


% --------------------------------------------------------------
% --------------------------------------------------------------
\begin{frame}[fragile]
  \frametitle{Rouletting}
  
  If a particle has too low of a weight, increase its weight or kill it.
  \vspace*{1 em}
  
  \begin{itemize}
  	\item if $w_i < w_{\min}$, roulette:
  	\begin{itemize}
  	    	\item Roulette Ratio, RR, $\equiv \frac{w_i}{w_{nom}}$
  	\item get $\xi$, a random number on $0 \leq \xi \leq 1$
  	\item if $\xi \geq$ RR:\\
  	      \hspace*{2 em} kill particle, $w_{i+1} = 0$ 
  	\item else:\\
  	      \hspace*{2 em} $w_{i+1} = w_{nom}$
  	\end{itemize}
  \end{itemize}

\end{frame}


% --------------------------------------------------------------
% --------------------------------------------------------------
\begin{frame}[fragile]
  \frametitle{Example}
  
  \begin{itemize}
  \item   A particle with $w_i = 0.9$ moves to a region with $w_{nom}$ is 3.5 and $w_{\min} = 1.5$
  
  \item We select $\xi = 0.23$
  
  \item What happens?
  \end{itemize}


\end{frame}


% --------------------------------------------------------------
% --------------------------------------------------------------
\section{\scshape Hybrid Methods}
%\subsection{Motivation}
\begin{frame}[fragile]
  \frametitle{Solving the TE}

\begin{columns}
  \begin{column}{0.5\textwidth}
  \begin{center}
  \underline{Monte Carlo}
  \end{center}
	\begin{itemize}
	\item Solution has associated statistical error
	\item Continuous phase space: ``gold standard answers"
	\item Can take a long time
	\item Good for streaming
	\item Optically thick = slow
	\end{itemize}
  \end{column}
  \begin{column}{0.5\textwidth}
  \begin{center}
  \underline{Deterministic}
  \end{center}
	\begin{itemize}
	\item Solution equally valid everywhere
	\item Discretized phase space: drives solution quality
	\item Can be fast
	\item Streaming = ray effects
	\item Good for optically thick
	\end{itemize}
  \end{column}
\end{columns}

\vspace*{1em}
\underline{Idea}: can we use deterministic and Monte Carlo methods together to lessen the weaknesses of each?

\vspace*{1em}
  $\rightarrow$ \textbf{Hybrid Methods}

\end{frame}


% --------------------------------------------------------------
%\subsection{CADIS}
\begin{frame}[fragile]
  \frametitle{Forward-Adjoint Relationship}
Define response with function $f(\ve{r}, E)$ in volume $V_d$ as
%
\begin{equation}
 R = \int_E \int_{V_f} f(\ve{r}, E) \phi(\ve{r}, E) dV dE 
 \label{eq:Response}
\end{equation}
%
\begin{columns}
  \begin{column}{0.5\textwidth}
	\begin{align}
  	H\phi &= q \quad \text{(forward)}\nonumber \\
  	%
  	H^{\dagger} \phi^{\dagger} &= q^{\dagger} \quad 
  	\text{(adjoint)}\nonumber
  	\end{align}
  \end{column}
  \begin{column}{0.5\textwidth}
  	\begin{align}
  	\langle H\phi, \phi^{\dagger} \rangle &= \langle H^{\dagger} \phi^{\dagger}, \phi \rangle \:, \text{and therefore} \nonumber \\
  	%
  	\langle q, \phi^{\dagger} \rangle &= \langle q^{\dagger}, \phi \rangle \nonumber
  	\end{align}
  \end{column}
\end{columns}
\vspace*{1 em}
If we let $q^{\dagger} = f(\ve{r}, E)$ then
%
\begin{equation}
 \langle q^{\dagger}, \phi \rangle = \langle f, \phi \rangle = R = \langle q, \phi^{\dagger} \rangle
 \label{eq:ResponseRedef}
\end{equation}
%
Eq.\ \eqref{eq:ResponseRedef} expresses that $\phi^{\dagger}$ represents the expected contribution of a source particle to the response given the source, $q$.

\end{frame}

% --------------------------------------------------------------
\begin{frame}[fragile]
  \frametitle{CADIS}
  
  \begin{enumerate}
  \item Define $q^{\dagger}$ as the local response of interest\\
  \item Coarse deterministic calculation to get $\phi^{\dagger}$ and $R$
  \end{enumerate}
% 
\begin{align}
  imp(\ve{r}, E) &= \frac{\phi^{\dagger}(\ve{r}, E)}{\langle q(\ve{r}, E), \phi^{\dagger}(\ve{r}, E) \rangle} = \frac{\phi^{\dagger}(\ve{r}, E)}{R} \\
  %
  \hat{q}(\ve{r}, E) &= \frac{\phi^{\dagger}(\ve{r}, E) q(\ve{r}, E)}{R} \\
  %
  w_0(\ve{r}, E) &= \frac{q(\ve{r}, E)}{\hat{q}(\ve{r}, E)} = \frac{R}{\phi^{\dagger}(\ve{r}, E)} 
  \label{eq:Importance}
\end{align}

Birth weights match weight targets, making this the \underline{C}onsistent \underline{A}djoint \underline{D}riven \underline{I}mportance \underline{S}ampling \underline{M}ethod

\end{frame}

% --------------------------------------------------------------
%\subsection{FW-CADIS}
\begin{frame}[fragile]
  \frametitle{FW-CADIS}

\begin{itemize}
\item We often what to optimize solutions in all of phase space\\
\item In this case the adjoint source needs to be a global forward solution: \underline{F}orward \underline{W}eighted-CADIS
\end{itemize}
%
\begin{columns}
  \begin{column}{0.5\textwidth}
  \begin{center}
  \textcolor{byellow}{To Optimize}
  \end{center}
	\begin{align}
  	&\phi(\ve{r}, E)\nonumber \\
  	%
  	\int&\phi(\ve{r}, E)\sigma_d(\ve{r}, E)\nonumber
  	\end{align}
  \end{column}
  %
  \begin{column}{0.5\textwidth}
  \begin{center}
  \textcolor{byellow}{Adjoint Source}
  \end{center}
  	\begin{align}
  	f(\ve{r}, E) &= \frac{1}{\phi(\ve{r}, E)}\nonumber \\
  	%
  	f(\ve{r}, E) &= \frac{\sigma_d(\ve{r}, E)}{\int\phi(\ve{r}, E)\sigma_d(\ve{r}, E)} \nonumber
  	\end{align}
  \end{column}
\end{columns}
\vspace*{1 em}
For example
%
\begin{equation}
 R = \int_E \int_{V_f} f(\ve{r}, E) \phi(\ve{r}, E) dV dE = \int_E \int_{V} \frac{1}{\phi(\ve{r}, E)} \phi(\ve{r}, E) dV dE \approx 1 \nonumber
\end{equation}

\end{frame}

% --------------------------------------------------------------
%\subsection{Challenges}
\begin{frame}[fragile]
  \frametitle{Example Impact}

	FW-CADIS works well for \textbf{most} deep penetration
	shielding problems
	%
	\begin{columns}
  	\begin{column}{0.5\textwidth}
  	\begin{figure}
  	\begin{center}
  		\includegraphics[height=2in,clip]{dlvn}
		\caption{Dog Legged Void Neutron Shielding Benchmark}
	\end{center}
  	\end{figure}
  	\end{column}
 	%
 	\begin{column}{0.5\textwidth}
 	\begin{figure}
  	\begin{center}
  		\includegraphics[height=2in,clip]{dlvn-lowVR}
  		\caption{MC 95\% CI RE using FW-CADIS, DLVN}
  	\end{center}
  	\end{figure}
  	\end{column}
	\end{columns}
  
\end{frame}




% --------------------------------------------------------------
% --------------------------------------------------------------
\section*{}
\begin{frame}[fragile]
  \frametitle{Questions?}
  \begin{center}
  \includegraphics[height=3in,clip]{../Class1/math07}  
  \end{center}
  
\end{frame}


\end{document}